\makeabstract
    {
    Development of IL-6 Neutralizing VHH Nanobodies to Prevent Cytokine Release Syndrome During CAR-T Cell Therapy
    }
    {
    Allie McLaughlin\textsuperscript{1,2}, Shahryar Khoshtinat Nikkhoi\textsuperscript{2}, Neil Savage\textsuperscript{2}, Jillian Baker\textsuperscript{2}, Anusuya Ramasubramanian\textsuperscript{2}, Eric Smith\textsuperscript{1}
    }
    {
    \textsuperscript{1} Dana-Farber Cancer Institute, Harvard Medical School, Boston, MA \\
\textsuperscript{2} Amherst College, Amherst, MA
    }
    {
    Chimeric antigen receptor (CAR) T-cell therapy has revolutionized the treatment of hematologic malignancies by recognizing and eliminating cancer cells. CAR T-cell therapy is frequently associated with cytokine release syndrome (CRS), a potentially life-threatening systemic inflammatory response caused by excessive immune activation. Interleukin-6 (IL-6) is a key mediator of CRS, and blockade of IL-6 signaling with tocilizumab is the current treatment. The large size of conventional monoclonal antibodies limits their incorporation into CAR T-cell constructs, highlighting the need for alternative therapeutic strategies.
    }
    {
    An initial library of 73 IL-6-specific VHH candidates was screened using phage display with human IL-6 for positive selection and human IL-11 for depletion. Candidate VHHs were expressed in E. coli TOP10F' cells, purified using nickel affinity chromatography, and normalized to 4000 nM. Functional IL-6 antagonism was assessed using a HEK293 IL-6 reporter assay, followed by bio-layer interferometry (BLI) to characterize binding kinetics. Lead candidates were evaluated by BLI epitope binning to determine whether they recognized distinct IL-6 epitopes, and sequence analysis was performed to identify complementarity-determining regions (CDRs). 
    }
    {
    Four VHH candidates (8, 1, 4, and 3) consistently reduced IL-6 reporter signaling to approximately 50\% of the IL-6-only control. BLI identified candidates 8 and 3 as the lead candidates, with binding affinities of 13.2 nM and 16 nM, respectively. Epitope binning demonstrated that candidates 3 and 8 recognize distinct IL-6 epitopes, supporting their potential use in a biparatopic VHH construct. Sequence analysis identified the CDRs of the lead candidates, providing a foundation for future VHH engineering.
    }
    {
    This work narrowed an initial pool of 73 IL-6-specific VHH candidates to two lead candidates, 3 and 8, with functional IL-6 antagonism and strong binding affinity. Their distinct IL-6 epitopes and characterized CDR sequences provide a foundation for next-generation VHH engineering, including development of a biparatopic construct and affinity maturation of candidate 8. These efforts may ultimately support the development of CAR T-cell therapeutics designed to mitigate CRS. 
    }

    \newpage
    